% Makra a styly obrazku prilohy C (Ukazkovy test, soubory src/18*.woo); prefix chU
\tikzset{
  chUlit/.style={v,minimum size=8.5mm,font=\small},
  chUin/.style={chUlit,fill=acc!25,draw=acc!80!black,very thick},
  chUcl/.style={v,minimum size=7.5mm,font=\scriptsize,fill=defc!8,draw=defc!80!black},
  chUa/.style={->,thick,draw=black!75,shorten >=1pt},
  chUc/.style={->,semithick,draw=defc!45,shorten >=1pt},
  chUcc/.style={->,semithick,draw=thmc!45,shorten >=1pt},
  chUabs/.style={->,very thick,draw=acc!80!black,shorten >=1pt},
}
% \chUpair{x}{y}{jmeno}{popisek kladneho}{popisek zaporneho}{styl kladneho}{styl zaporneho}
% promennovy gadget: dva vrcholy (jmeno)p a (jmeno)n spojene hranami v obou smerech
\newcommand{\chUpair}[7]{%
  \node[#6] (#3p) at (#1,#2) {#4};
  \node[#7] (#3n) at ({#1+1.5},#2) {#5};
  \draw[chUa] (#3p) to[bend left=28] (#3n);
  \draw[chUa] (#3n) to[bend left=28] (#3p);}
% \chUtri{x}{y}{jmeno}{index klauzule}: orientovany trojuhelnik c^1 -> c^2 -> c^3 -> c^1 se stredem (x,y);
% vrcholy se jmenuji (jmeno)a, (jmeno)b, (jmeno)c
\newcommand{\chUtri}[4]{%
  \node[chUcl] (#3a) at (#1,{#2+0.85}) {$c_{#4}^1$};
  \node[chUcl] (#3b) at ({#1+1.1},{#2-0.6}) {$c_{#4}^2$};
  \node[chUcl] (#3c) at ({#1-1.1},{#2-0.6}) {$c_{#4}^3$};
  \draw[chUa] (#3a)--(#3b);
  \draw[chUa] (#3b)--(#3c);
  \draw[chUa] (#3c)--(#3a);}
